\section{Masked-cause inference: brief primer}
\label{sec:background}

We summarize the masked-cause identifiability framework
\citep{towell2026masked,heitjan1991ignorability,gill1997coarsening}.
Readers familiar with that work, or with the companion applications
\citep{towell2026scrnacoarsening,towell2026spatialcoarsening}, can
skim to \cref{sec:translation}.

Consider $m$ independent components in series. Each component $j$ has
a parametric distribution $f_j(\cdot;\theta_j)$ for its lifetime
$T_{ij}$ in system $i$. The system fails at $T_i = \min_j T_{ij}$
with cause $K_i = \arg\min_j T_{ij}$. We observe $T_i$ but not $K_i$
directly; instead we observe a \emph{candidate set} $C_i \subseteq
\{1,\ldots,m\}$ thought to contain the failed component.

The joint likelihood factorizes as
\begin{equation}
  f(t, k, c; \theta)
  = h_k(t; \theta_k) \prod_{\ell} R_\ell(t; \theta_\ell)
    \cdot \Prob_\theta\{C_i = c \mid T_i = t, K_i = k\},
  \label{eq:bg-joint}
\end{equation}
where $h_\ell$ is the component-$\ell$ hazard and $R_\ell$ its
survival function. The third factor is the unknown coarsening (or
masking) mechanism.

Three coarsening-at-random conditions allow the masking factor to
drop from the score:

\begin{condition}[Support, C1]
\label{cond:c1}
$\Prob\{K_i \in C_i\} = 1$: the true cause lies in the candidate set.
\end{condition}

\begin{condition}[Symmetry, C2]
\label{cond:c2}
$\Prob_\theta\{C_i=c \mid T_i=t, K_i=j\} = \Prob_\theta\{C_i=c \mid
T_i=t, K_i=k\}$ for all $j,k \in c$: the coarsening does not depend on
which element of the candidate set is the true cause.
\end{condition}

\begin{condition}[Parameter independence, C3]
\label{cond:c3}
$\Prob_\theta\{C_i=c \mid T_i, K_i\}$ does not depend on $\theta$: the
coarsening parameters are distinct from the component parameters.
\end{condition}

Under all three, the likelihood collapses to
\begin{equation}
  L_i(\theta) \propto \left[\prod_\ell R_\ell(t_i;\theta_\ell)\right]
  \cdot \left[\sum_{j \in c_i} h_j(t_i; \theta_j)\right],
  \label{eq:bg-c1c2c3}
\end{equation}
that is, the coarsening mechanism becomes ignorable for likelihood
inference.

\citet[Theorem 8]{towell2026masked} establish identifiability under
these conditions:

\begin{theorem}[Identifiability under C1--C2--C3]
\label{thm:bg-id}
Under \cref{cond:c1,cond:c2,cond:c3} and standard regularity
(in particular, positive probability assigned to both exact-event
and censored-type observations), full column rank $m$ of the
augmented candidate-set matrix $\tilde{C} \in \{0,1\}^{n \times m}$
is \emph{sufficient} for $\theta$ to be identifiable. Separability
is necessary but not in general sufficient (a separating mechanism
can leave $\tilde{C}$ rank-deficient, e.g. $m=4$ with candidate
sets $\{1,2\},\{3,4\},\{1,3\},\{2,4\}$).
\end{theorem}

\emph{Singleton} candidate sets ($|c_i|=1$) are the critical
ingredient: each isolates a single component contribution in
\eqref{eq:bg-c1c2c3} and adds an identifying unit-vector row to
$\tilde{C}$. When broader candidate sets are confounded (the matrix
$\tilde{C}$ is rank-deficient), singletons restore the missing rank.

This framework has been ported to single-cell RNA sequencing zero
inflation \citep{towell2026scrnacoarsening} and to
spatial-transcriptomics deconvolution
\citep{towell2026spatialcoarsening}. We now port it to programmatic
weak supervision. The candidate-set structure here is combinatorial:
each labeling function emits a vote (or abstains), and the
intersection of the votes defines the candidate set over labels. The
coarsening mechanism is the joint LF noise model, and the three
conditions become statements about how LF errors relate to the true
label and to the label-model parameters.
